% Nouveau EGA V (cote 25): archive cote 25, site page 32 (PDF page 33).
% typescript. Draft compared: none.
% AI-transcribed by Claude Opus 5.5 (claude-opus-5-5, adaptive thinking, effort high), Message Batches API, batch ega25b (msgbatch_01A7n5YuXubUPNK1zsepi8RG), 2026-11-13; not yet checked by eye.
\scanpage{25/32}
\mspage{28}

8.2. Soit Z l'ensemble des $\hand{\xi\in}\check{P}$ tels que dim $Y_{\hand{\xi}}$ $\hand{>}$ dim X - 1,\hand{\unsure{\xi}}
et supposons que pour toute composante irr $X_i$ de X, on ait dim $f(X_i)$ $\hand{>}$ 0;
alors Z est de codimension 2 dans $\check{P}$. \xout{(ixexxdexdimensionxxxx=2)} Cela
\inserted[typed]{2.1. et 2.2.}
résulte de \xout{\unsure{2x5xxxx2x3xxx(quiximpliquentxquexxlesxXx$X_i$xsontxirréductibles}}
\xout{\unsure{etxdominentxP}} (qui implique que toute comp irr de $\underline{Y}$ domine $\check{P}$ ) et de
la théorie de la dimension pour le morphisme $\underline{Y} \hand{\to} \check{P}$ . A partir de ce
résultat \xout{\unsure{(quixxxsuffitxde}} on peut donner comme corollaire le cas où on
part d'une partie \inserted[typed]{fermée} \xout{\illegible{1}} Z de X et qu'on r egarde la dimension des
images inverses $Z_{\hand{\xi}}$ dans les $Y_{\hand{\xi}}$ ( $\hand{\xi\in}\check{P}$ ), et on peut même prendre pour
Z l'ensemble des $\hand{\xi\in}\check{P}$ tels que il existe une co mposante irr. de $T_{k(\hand{\xi})}$
dont la trace sur $Y_{\hand{\xi}}$ soit de dimension trop grande (NB on suppose toujour
que pour toute comp irréd $Z_i$ de Z, on ait dim $f(Z_{\hand{i}})$ $\hand{>}$ 0). Enfin, l'énoncé
le plus précis dans cette direction, et qui résulte aisément du premier
énoncé (pour X irréductible) et de 2.7. , est l'énoncé modulé suivant:
F étant cohérent sur X, \xout{soit} supposons que pour tout cycle premier associ
T pour F, on ait dim f(T)$\hand{>}$ 0\unsure{;} alors l'ensemble des $\hand{\xi\in}\check{P}$ tels que
$\hand{\phi_\xi}$ \inserted{ne} soit pas $F_{k(\hand{\xi})}$-régulier est (constructible et) de codimension$\hand{\geqslant}$ 2.
\margin{\textbar} (La notation pour $\hand{\phi_\xi}$ est celle du n°5). On peut donner ceci comme
\margin{\textbar} \inserted[typed]{proposition principale}\xout{\unsure{Dexlaxmêmexfaçon}}, et énoncer les \xout{autres} assertions qui précédaient
\margin{\textbar} comme corollaires, la démonstration procédant via l'un des corollaires.

\deleted{Autre corollaire: Soient T,T' des parties femmées de X, et supposons
qu'on \xout{mait} \xout{\unsure{pasxdimxexdimxim}} T / $T'_{\unsure{\xi}}$ ( étant le point générique de
\inserted[typed]{il existe une co mp irr $T_i$ de T telle que}
P ), i.e. $f(T_i)$ est infini et $T_i$ / T' (cf 2.4.); alors l'ensemble Z des
$\check{P}$ tels que T T' est de codimension 2. Pour le voir, on se
ramène au cas T' T (remplaçant T' par T T'), et où T est irréductible
(le remplaçant par une composante irréductible convenable). Alors, \xout{xx}
d'après \inserted[typed]{$Z_o$} ce qu'on a vu, si $\tilde{P}$ n'appartient pas à une certaine partie
fermée de \unsure{codimension} 2, chacun\xout{e} des \xout{\unsure{fibres}} fermés T', T}

\mspage{32}
